\section{La partícula como clase de secciones persistentes}
\label{vi:sec:particula}

La construcción material comienza en el estado que producen APP, TRIT y
TPK. APP proporciona las posiciones y las dos operaciones aritméticas;
TRIT distingue el régimen y la orientación de sus transiciones; TPK
transporta esa información junto con los cocientes, las incidencias de
frontera y la memoria. Al prolongar el estado, una evaluación finita
permanece vinculada a las evaluaciones que la preceden. La noción de
partícula se apoya en esta compatibilidad y en su representación material.
Su masa será un observable de la sección así construida.

Esta secuencia determina el significado de los tres objetos que se
utilizarán: el estado genealógico conserva la extensión completa; la ruta
registra su historia observable; la representación material hace actuar
esa historia sobre un espacio de estados y permite evaluar observables.
La distinción resulta esencial incluso cuando los tres objetos comparten
una misma etiqueta electrónica o una misma coordenada del atlas.

\subsection{Extensión, historia observable y sección compatible}
\label{vi:subsec:extension-persistente}

Sean $\mathcal S_m$ los estados admisibles supervivientes hasta profundidad
$m$, obtenidos por la prolongación del TPK, y sean
$\rho_{n,m}:\mathcal S_n\to\mathcal S_m$, para $n\geq m$, sus
aplicaciones de restricción. La compatibilidad del refinamiento significa
\begin{equation}
 \rho_{m,m}=I,\qquad
 \rho_{n,m}\rho_{p,n}=\rho_{p,m}\quad(p\geq n\geq m).
 \label{vi:eq:sistema-estados}
\end{equation}
Una extensión superviviente es un elemento
\begin{equation}
 x=(x_m)_{m\geq0}\in\Omega^{\mathrm{surv}}_{\mathrm{HMT}}
 :=\varprojlim_m\mathcal S_m,
 \qquad \rho_{n,m}(x_n)=x_m.
 \label{vi:eq:extension-superviviente}
\end{equation}
La existencia y la admisibilidad de estos estados proceden de la
construcción del núcleo formal. La expresión de límite inverso describe
las compatibilidades que conserva un estado ya generado.

Sea $\mathcal R_m$ el espacio de historias observables a profundidad $m$.
El lector de ruta $\operatorname{sh}_m:\mathcal S_m\to\mathcal R_m$
y la restricción $\operatorname{tr}_{n,m}$ satisfacen
\begin{equation}
 \operatorname{tr}_{n,m}\operatorname{sh}_n
 =\operatorname{sh}_m\rho_{n,m}.
 \label{vi:eq:naturalidad-ruta}
\end{equation}
Así, el lector transmite las restricciones del estado a sus historias.
La ruta de $x$ es la familia
$\gamma_x=(\gamma_x^{[m]})_m$, donde
$\gamma_x^{[m]}=\operatorname{sh}_m(x_m)$.

\begin{proposition}[Compatibilidad de la historia publicada]
\label{vi:prop:historia-compatible}
Para toda extensión \eqref{vi:eq:extension-superviviente}, sus historias
finitas cumplen
$\operatorname{tr}_{n,m}(\gamma_x^{[n]})=\gamma_x^{[m]}$.
\end{proposition}
\begin{proof}
Aplicando \eqref{vi:eq:naturalidad-ruta} al estado $x_n$ se obtiene
\[
 \operatorname{tr}_{n,m}(\gamma_x^{[n]})
 =\operatorname{sh}_m(\rho_{n,m}x_n)
 =\operatorname{sh}_m(x_m)=\gamma_x^{[m]}.
\]
La segunda igualdad utiliza precisamente la condición de límite inverso.
\end{proof}

La proposición establece el descenso de la extensión a la historia.
Su contenido no incluye la inyectividad del lector: dos extensiones
pueden publicar una misma historia finita y conservar memorias distintas.
Por ello, al reutilizar una ruta se mantiene su estado de procedencia,
en lugar de reconstruirlo a partir de una cifra o de una palabra visible.

Sobre cada historia $\gamma_x^{[m]}$ se considera el fibrado discreto
$\mathcal B_{\gamma_x^{[m]}}$ de estados materiales de la representación
declarada. Sean
$\mathscr E_m(x)=\Gamma(\mathcal B_{\gamma_x^{[m]}})$ sus espacios
de secciones, con restricciones
$p_{n,m}:\mathscr E_n(x)\to\mathscr E_m(x)$ compatibles con las
del estado. Una sección persistente es una familia
\begin{equation}
 s=(s_m)_{m\geq m_0},\qquad s_m\in\mathscr E_m(x),\qquad
 p_{n,m}s_n=s_m\quad(n\geq m\geq m_0).
 \label{vi:eq:seccion-persistente}
\end{equation}
La profundidad inicial $m_0$ permite describir un estado desde cualquier
corte a partir del cual su representación esté definida. La información
anterior que intervenga en sus observables permanece en el registro
transportado por la extensión.

\subsection{Identidad bajo refinamiento y equivalencia de representaciones}
\label{vi:subsec:identidad-material}

Fijemos los grupos de simetrías internas $G_m$ que actúan en las fibras
materiales y preservan los observables declarados. Dos secciones
persistentes $s$ y $s'$ representan la misma identidad material cuando,
a partir de algún nivel $m_1$, existe una familia $g_m\in G_m$ tal que
\begin{equation}
 s'_m=g_ms_m,\qquad
 p_{n,m}g_n=g_mp_{n,m},
 \label{vi:eq:equivalencia-persistente}
\end{equation}
y coinciden los caracteres observables del registro genealógico. Las
simetrías admisibles se fijan antes de formar el cociente; su acción
transporta la hoja, la orientación, la memoria y la representación con
las leyes correspondientes.

\begin{proposition}[Equivalencia compatible de las secciones]
\label{vi:prop:equivalencia-persistente}
La relación \eqref{vi:eq:equivalencia-persistente}, sobre familias que
portan los mismos tipos de representación y observables, es una relación
de equivalencia. Sus clases son independientes de retirar un número
finito de niveles iniciales, siempre que se conserve su información
genealógica transportada.
\end{proposition}
\begin{proof}
La identidad de cada $G_m$ da reflexividad. La inversa de una familia
compatible vuelve a ser compatible: de
$p_{n,m}g_n=g_mp_{n,m}$ se deduce
$p_{n,m}g_n^{-1}=g_m^{-1}p_{n,m}$. Esto da simetría.
Si $s'=gs$ y $s''=hs'$, a partir del mayor de los dos niveles iniciales
se tiene $s''_m=h_mg_ms_m$ y
$p_{n,m}h_ng_n=h_mg_mp_{n,m}$. La composición preserva también los
observables, y prueba transitividad. Todas estas relaciones se formulan
sobre una cola común; cambiar el punto inicial dentro de esa cola
conserva la clase. La última afirmación se refiere al cambio de
presentación, no a borrar memoria del estado.
\end{proof}

\begin{definition}[Partícula material persistente]
\label{vi:def:particula}
Una partícula HMT--MD es una clase $[s]_{\mathrm{pers}}$ de secciones
\eqref{vi:eq:seccion-persistente} bajo
\eqref{vi:eq:equivalencia-persistente}, acompañada de su extensión y
registro de procedencia, que admite una representación material
irreducible o un bloque mínimo invariante y una salida espectral estable
bajo el acoplamiento especificado.
\end{definition}

La definición reúne condiciones de distintos niveles. La prolongación
y el registro pertenecen al estado HMT. La representación y el
acoplamiento pertenecen a su realización material. La estabilidad
espectral se predica del operador y del dominio utilizados por ese
acoplamiento. Esta organización permite estudiar conjuntamente identidad,
carga, espín y masa sin convertir un observable particular en definición
de todos los demás.

El retorno nonádico ilustra la distinción entre identidad y estado
instantáneo. Para el lector de fase $g$ y el registro de memoria
$\mathcal B$ pueden coexistir
\begin{equation}
 g(k+9)=g(k),\qquad
 p_{k+9,k}(s_{k+9})=s_k,\qquad
 \mathcal B_{k+9}\ne\mathcal B_k.
 \label{vi:eq:persistencia-fase-memoria}
\end{equation}
La primera igualdad es periódica; la segunda expresa persistencia; la
tercera registra el avance. El retorno de fase pertenece a la historia
de una misma identidad material, mientras el estado conserva que se ha
efectuado una vuelta adicional.

\subsection{La especie como fibra: alcance del atlas finito}
\label{vi:subsec:especie-fibra}

El atlas construido por las operaciones APP--TRIT--TPK tiene soporte
$\mathcal C_{81}=\{1,\ldots,9\}^2$. Sus cuatro orientaciones por
celda forman el dominio de $324$ rutas orientadas. La aplicación de
registro $\mathcal R$ distingue $56$ familias de ruta, mientras el
clasificador interno
\begin{equation}
 q:\mathcal C_{81}\longrightarrow\Sigma_{13}
 \label{vi:eq:clasificador-especies}
\end{equation}
reúne $13$ tipos de familia y nivel. Estos censos y sus reglas se
desarrollan en el cuerpo dedicado al atlas; aquí se precisa qué objeto
transmiten a la definición de partícula.

Para $y\in\Sigma_{13}$, la fibra $q^{-1}(y)$ conserva todas sus
celdas y registros. Fijada la ordenación interna de cada fibra de
$\mathcal R$, sea $k(c)$ el rango de la celda dentro de ella. Su
presentación finita es
\begin{equation}
 \mathfrak S_y^{\mathrm{atlas}}
 =\{(\mathcal R(c),k(c)):q(c)=y\}.
 \label{vi:eq:multiseccion-atlas}
\end{equation}
El rango se añade para separar celdas que poseen el mismo registro de
familia; no selecciona una masa. El par $(\mathcal R(c),k(c))$
recupera la celda dentro de la carta ordenada porque el segundo dato
individualiza su posición en la fibra del primero.

Si $\pi_{81}:\Omega^{\mathrm{surv}}_{\mathrm{HMT}}\to
\mathcal C_{81}$ es el lector celular, la elevación de la especie es
\begin{equation}
 \widetilde{\mathfrak S}_y
 =\{x\in\Omega^{\mathrm{surv}}_{\mathrm{HMT}}:
                  q(\pi_{81}x)=y\}.
 \label{vi:eq:multiseccion-enriquecida}
\end{equation}
Esta preimagen contiene extensiones, mientras
\eqref{vi:eq:multiseccion-atlas} contiene presentaciones finitas. La
construcción y el dominio del lector $\pi_{81}$ forman parte del paso
entre ambas. Una partícula concreta se realiza como sección persistente
dentro de la multisección enriquecida correspondiente.

\begin{theorem}[Punto fijo central y selección equivariante]
\label{vi:thm:selector-central}
Sea $\rho_c(r,s)=(10-r,10-s)$ la inversión celular. En las fibras
$\rho_c$-estables del clasificador \eqref{vi:eq:clasificador-especies},
con acción trivial sobre la etiqueta $y$, una selección puntual
equivariante sólo puede tomar valores en la celda $(5,5)$. La fibra
electrónica de nivel cero contiene ese único punto fijo; ninguna otra
fibra admite tal selección.
\end{theorem}
\begin{proof}
Resolver $\rho_c(r,s)=(r,s)$ da $r=s=5$. Si $a(y)\in q^{-1}(y)$
es una selección equivariante y $y$ queda fijo, entonces
$a(y)=a(\rho_c y)=\rho_c a(y)$. Su valor debe ser el punto fijo
$(5,5)$. El atlas sitúa esa celda en la fibra electrónica de nivel
cero, cuyo cardinal es uno. Para cualquier otra fibra, la condición
de selección produciría un punto fijo que no pertenece a ella.
\end{proof}

La consecuencia es constructiva: las demás especies se presentan por
órbitas o multisecciones, y una realización que necesite una línea
particular debe declarar su representación y su acoplamiento. El
electrón proporciona el caso central de rango uno sin convertir todas
las familias en copias de ese caso. Además, la celda $(5,5)$, la firma
central cruda, la firma regional $555555$ y el origen relativo de una
recta de acción conservan sus dominios propios. La unicidad del punto
no identifica esos registros entre sí.

\subsection{Del registro material al observable de masa}
\label{vi:subsec:particula-observable}

La cadena utilizada por el observable de masa conserva el orden
\begin{equation}
 \text{extensión}\longrightarrow\text{ruta}\longrightarrow
 \text{registro}\longrightarrow\text{firma entera}\longrightarrow
 \text{carácter}\longrightarrow\text{torres}\longrightarrow
 \text{operador material}.
 \label{vi:eq:cadena-material}
\end{equation}
Cada flecha retiene los datos necesarios para la siguiente. La firma es
una lectura del registro; el carácter evalúa esa firma con los
coeficientes internos ya construidos; las torres incorporan su
refinamiento; el operador actúa finalmente en la representación de la
fibra seleccionada por el acoplamiento. Su salida puede ser escalar,
un multiplete o una distribución espectral, según esa representación.
La comparación metrológica tiene lugar después de esta construcción.

Una especie queda descrita por carga, representación de espín, color,
sabor, generación, familia de ruta y operador material. Estos datos
explican por qué dos preparaciones con una masa coincidente pueden
conservar distinta información, y por qué una misma identidad puede
publicar observables diferentes bajo acoplamientos distintos. La
identidad reside en la clase persistente y en sus transportes admisibles;
la medida es una lectura de esa identidad en el codominio declarado.

La misma jerarquía separa una extensión persistente de una sección de
horizonte finito. Si $\operatorname{Ext}_L(s)\ne\varnothing$ y
$\operatorname{Ext}_{L+1}(s)=\varnothing$, el horizonte $L$ registra
una obstrucción de prolongación. Una participación intermedia de este
tipo en una composición no crea, por sí misma, una nueva especie
asintótica. Análogamente, una resonancia exige el operador de
supervivencia y su polo, además de la identidad de la sección que
transporta. Se preserva así la diferencia entre persistencia,
prolongación finita y evaluación espectral.
